% Source: Paper_A_新版完整论文提纲_vFinal_Draft.md §1
\section{引言}

\subsection{研究背景}

全景房间布局标注需要工人同时处理：

\begin{itemize}
  \item 相机所在空间的识别；
  \item Scope/OOS 判断；
  \item 房间边界和角点几何；
  \item 模型初始化的识别与纠正；
  \item 遮挡、开放边界和邻接空间；
  \item 任务过程与完整性。
\end{itemize}

Semi-Auto 初始化可能降低操作负担，但也可能诱发：

\begin{itemize}
  \item blind trust；
  \item correction failure；
  \item over-correction；
  \item undercoverage；
  \item 结构性无效提交。
\end{itemize}

不同工人在一般任务、高风险任务和模型错误条件下的表现并不一定相同。单一的全局工人排名可能无法充分利用这种差异，但高维个人画像又容易受到小样本和支持不足影响。

\subsection{研究缺口}

现有方案通常存在以下不足：

\begin{enumerate}
  \item 把 PreScreen 仅视为准入，而没有检验其跨阶段预测价值；
  \item 将工人与同行一致性等同于外部正确性；
  \item 忽略结构性无效提交导致的选择性可计算问题；
  \item 在不同难度任务分发下直接按原始平均质量排名工人；
  \item 在小样本下无条件发布个体条件画像；
  \item 将路由简化为一个评分公式，而未定义 offer、容量、追加、聚合和终态；
  \item 只在 resolved 子集评价质量，忽视 unresolved 的选择性分母；
  \item 使用平衡试验分布替代真实生产任务分布。
\end{enumerate}

\subsection{方法路线}

本文采用：

\begin{verbatim}
高信息 P1
→ C1 基础测量
→ C2-B 共同桥接与风险确认
→ C2-A-RP 精度补测
→ T1 Semi 条件效应
→ Strong Global 对比 Support-aware Full V1
\end{verbatim}

并始终分离：

\begin{verbatim}
公开/纠错 GT 外部正确性
Geometry LOO 同行一致性
worker-caused structural failure
模型辅助风险
Calibration 支持范围
\end{verbatim}

\subsection{一级贡献}

\subsubsection{贡献一：三轴工人测量}

本文分别估计：

\begin{itemize}
  \item $Q^{GT}_u$：相对 operational GT 的外部正确性；
  \item $R^{peer}_u$：与独立同行的几何兼容性；
  \item $F^{struct}_u$：可归责于工人的结构失败倾向。
\end{itemize}

三者不压缩为单一总可靠度；P1 预测链和多峰共识作为次级效度证据。

\subsubsection{贡献二：双波 Calibration 与部分收缩}

C1 提供基础外部质量、同行结构和任务差异；C2-B 通过 common anchor 与 diverse
bridge 恢复横向可比性并确认风险差异，C2-A-RP 只补关键估计精度。
小样本工人使用部分收缩，不发布不稳定的高维画像。

\subsubsection{贡献三：支持度感知的前瞻政策试验}

V1 比较 Strong Global 与 Support-aware Full。Full 只能使用获得跨阶段支持的条件分量，
并在支持不足时局部或整体退化为 Global。两臂共享容量、替补、动态追加和 GT-blind 聚合条件；
完整工程合同放入方法附录。

P1 跨阶段预测、多峰共识、GT 冲突解释和反例库是次级创新或可信度基础设施，
不是替代三项核心贡献的主叙事。

\subsection{研究问题}

\subsubsection{RQ1：模型辅助效应}

Semi-Auto 是否在不降低交付质量的前提下降低 active time，其效应是否随任务风险和工人纠错能力变化？
确认性层级为 delivery-adjusted quality 非劣、active time 优势和
\texttt{mode} $\times$ \texttt{risk\_assist} 交互；blind trust、correction failure、
over-correction 和 Model Issue recognition 为机制性次级结果。

\subsubsection{RQ2：画像测量和跨阶段效度}

P1 诊断和 C1/C2 三轴工人状态是否具有跨阶段预测价值，哪些有限的 failure-family
component 能获得可复现支持？RQ2 以效应量、方向、支持范围和预测性能为主，
不要求每阶段分别达到传统显著性。

\subsubsection{RQ3：路由增量价值}

在相同工人池、容量、替补、追加和聚合规则下，Support-aware Full 是否相对
Strong Global 改善交付调整质量、非交付率或标注成本？V2 外部支持审计是可选扩展，
不作为第四个主要研究问题。
